\maketitle
\tableofcontents

\chapter*{Dok.~310 -- Die Resonanzgeometrie als übergeordnete Referenz}
\addcontentsline{toc}{chapter}{Dok.~310 -- Die Resonanzgeometrie als übergeordnete Referenz}

\section*{Worum es geht}

Der wiederkehrende Vorwurf gegen FFGFT lautet: „$\xi$ ist an die
Leptonmassen angepasst, also zirkulär." Diese Notiz zeigt, dass der
Vorwurf auf einem Missverständnis der Hierarchie beruht.

Die richtige Ordnung ist: \textbf{übergeordnet} ist die Geometrie, die
Resonanzen besitzt — dimensionslos, legt alle Verhältnisse fest.
Der \textbf{Anker} ist ein einziger Messwert — trägt keine Struktur,
nur die Skala. Der Messwert ist der Kammerton, nicht die Musik.

Nicht Teil der A-Serie.

\section*{Konventionen}

\noindent\textbf{Zeichenerklärung.}

\medskip
\begin{tabular}{@{}lp{11cm}}
\textbf{[K]}       & \emph{Kernableitung} — aus den Grundsetzungen
                     hergeleitet, reproduzierbar. \\[4pt]
\textbf{[S]}       & \emph{Strukturaussage} — begründet, nicht vollständig
                     bewiesen. \\[4pt]
\textbf{[SETZUNG]} & \emph{Setzung} — deklarierte Annahme. \\
\end{tabular}

\medskip
\noindent\textbf{Zentrale Größen.}

\medskip
\begin{tabular}{@{}lp{10.5cm}}
$Q$ & Koide-Verhältnis $(\sum m)/(\sum\sqrt m)^2$; misst $2/3$ [K]. \\[4pt]
$\xi = 4/30000$ & Fundamentalparameter; Form [K], Magnitude $5^4$ [SETZUNG, P33]. \\[4pt]
$r_\ell, p_\ell$ & Vorfaktoren und Exponenten der Massenleiter (Dok.~190). \\[4pt]
$v$ & Vakuumerwartungswert; setzt die absolute Skala. \\[4pt]
$S^2 = 4\pi,\ S^3 = 2\pi^2$ & Sphären-Oberflächen der T$^4$-Randgeometrie. \\[4pt]
$\mu_0,\varepsilon_0$ & Vakuumkonstanten; folgen aus $\alpha,e,c$ (Dok.~013). \\
\end{tabular}

\chapter{Die drei Ebenen}

\begin{center}
\begin{tabular}{lll}
\toprule
Ebene & Inhalt & Status \\
\midrule
0 — Messung & $m_\mu/m_e = 206{,}77$, $m_\tau/m_e = 3477$ & gemessen \\
1 — Relation & $Q = 2/3$ auf $10^{-5}$ & parameterfrei [K] \\
2 — Kodierung & $m_\ell = r_\ell\,\xi^{p_\ell}\,v$ & Schreibweise \\
\bottomrule
\end{tabular}
\end{center}

Der einzige Fehler wäre, Ebene 2 als Ebene 0 zu verkaufen — also zu
behaupten, $\xi$ sei der Messwert-Ausgang. Solange der Status jeder
Ebene deklariert ist, ist die Darstellung ehrlich und durchsichtig
zugleich.

\chapter{Was die Geometrie vorgibt}

Die Resonanzstruktur macht drei vollständig dimensionslose Aussagen:

\textbf{(1) Koide-Relation [K].} Die drei $\sqrt m$ bilden einen Winkel
von $44{,}9997^\circ$ zum Vektor $(1,1,1)$:
\[
  Q = \frac{m_e + m_\mu + m_\tau}
           {(\sqrt{m_e}+\sqrt{m_\mu}+\sqrt{m_\tau})^2} = \frac{2}{3}.
\]
Kein Messwert nötig — nur die Verhältnisse der Massen zueinander.

\textbf{(2) Generationsleiter [K].} Die geladenen Leptonen realisieren
die Exponenten $p_e = 3/2$, $p_\mu = 1$, $p_\tau = 2/3$ — eine
\textbf{geometrische Folge} mit konstantem Verhältnis
$q = 2/3$:
\[
  p_{n+1} = \tfrac{2}{3}\,p_n, \qquad
  \tfrac{3}{2} \xrightarrow{\times 2/3} 1 \xrightarrow{\times 2/3} \tfrac{2}{3}.
\]
Das Verhältnis $q = 2/3$ ist derselbe Wert wie der Koide-Wert
$Q = 2/3$ — dieselbe Zahl ordnet die Leiter und die Massenrelation
(Dok.~258). Die Folge ist an den drei gemessenen Massen nachprüfbar;
dass andere Sektoren auf anderen Sprossen sitzen, berührt das nicht.

\textbf{(3) Harmonische Vorfaktoren [K].} Die Vorfaktoren
$r_e = 4/3$, $r_\mu = 16/5$, $r_\tau = 25/9 = (5/3)^2$ — Intervalle im
Euler-Tonnetz.

Diese Aussagen legen \textbf{alle} Verhältnisse fest — ohne einen
einzigen Messwert (\texttt{ffgft\_310\_massenleiter.py}):
\[
  m_\mu/m_e = 207{,}85 \ (+0{,}52\%), \qquad
  m_\tau/m_e = 3531{,}6 \ (+1{,}56\%).
\]

\chapter{Wo der Messwert eintritt}

Erst für die \textbf{absolute Skala} braucht man einen Anker:
$m_e = 0{,}511$~MeV. Er sagt nicht, wie sich die Massen verhalten —
das sagt die Geometrie. Er sagt nur: „diese Struktur, in kg gemessen,
fängt hier an." Man könnte genauso $m_\mu$ oder $m_\tau$ nehmen — die
Verhältnisse blieben identisch.

\textbf{Die Musik-Analogie.}

\begin{center}
\begin{tabular}{ll}
\toprule
Musik & Leptonen \\
\midrule
Intervalle ($2{:}1$, $3{:}2$, $5{:}4$) & $Q=2/3$, $1/3$-Leiter, harm.\ $r$ \\
Kammerton $a' = 440$~Hz & $m_e = 0{,}511$~MeV \\
\bottomrule
\end{tabular}
\end{center}

Man kann $a' = 440$ oder $442$~Hz wählen — die Musik bleibt dieselbe.
Der Anker ist Konvention, die Resonanzstruktur ist das Übergeordnete.

\chapter{Warum das die Zirkularität auflöst}

Der Vorwurf: „$\xi$ ist an die Massen angepasst." Die Antwort: Ja — die
absolute \emph{Skala} (der Kammerton) ist gemessen. Aber die
\emph{Struktur} (die Resonanzen) ist nicht angepasst. Sie wird durch
$Q = 2/3$ \textbf{bestätigt}, nicht eingestellt.

Der Beweis steht in der Restabweichung
(\texttt{ffgft\_310\_koide\_referenz.py}):
\[
  Q_\text{gemessen} - \tfrac{2}{3} = -6{,}2 \times 10^{-6} \neq 0.
\]
Ein Fit hätte $Q$ exakt getroffen ($0{,}000\%$). Die Geometrie trifft
$Q$ auf $10^{-5}$ \emph{mit Rest} — die Signatur einer echten
Vorhersage.

\chapter{Die Wahl des Ausgangspunkts ist Konvention}

Es gibt \textbf{eine} dimensionslose Zahlenwelt. $\xi$, $\alpha$, $Q$,
$m_\mu/m_e$ sind verschiedene Koordinaten darin — keine ist
„fundamentaler". Die Zirkularität ist nicht $\xi$-spezifisch: „warum
$\alpha = 1/137$?", „warum $Q = 2/3$?", „warum $\xi = 1/7500$?" sind
alle offen. Jede Theorie muss irgendwo eine dimensionslose Zahl setzen.

Der einzige echte Unterschied ist \textbf{eine vernetzte Setzung vs.
viele isolierte}: das SM hat 3 unabhängige Leptonmassen; FFGFT spart
über $Q = 2/3$ eine Zahl und verbindet über $\xi$ den Leptonsektor mit
der Kosmologie ($a_0$, $H_0$). Der Gewinn ist real, aber bescheiden —
kein „alles aus Nichts", sondern eine Reduktion mit Richtung.

\chapter{Der Entdeckungsweg und die Geometrie der \texorpdfstring{$\pi$}{pi}-Faktoren}

Der Korpus stellt $\xi \to$ Massen $\to$ Kosmologie dar (aufgeräumt).
Der tatsächliche Entdeckungsweg war: (1)~QM/RT-Unverträglichkeit als
Motiv, (2)~die Higgs-Formel als erster Ausdruck, (3)~Vakuum und seine
Konstanten, (4)~erst danach Massen, Koide, Tonnetz.

\section{Higgs-Formel als zweiter unabhängiger Zugang [K]}

\[
  \xi = \frac{m_h^2}{64\pi^3\, v^2} = 1{,}30 \times 10^{-4}
\]
trifft $\xi = 4/30000$ auf $2\%$, aus völlig anderen Eingaben
(Higgs-Masse, Vakuumwert) als die Leptonmassen
(\texttt{ffgft\_310\_higgs\_zugang.py}). Zwei unabhängige Sektoren,
ein $\xi$ — das war das Signal.

\section{$\pi^3$ ist Geometrie, keine Schleife [K]}

Der Faktor $\pi^3$ ist \emph{kein} Quanten-Schleifenfaktor, sondern das
Produkt der Sphären-Oberflächen der kompakten T$^4$-Randgeometrie
(\texttt{ffgft\_310\_pi\_geometrie.py}):
\[
  16\pi^3 = (4\pi)\cdot(4\pi^2) = S^2 \cdot 2S^3.
\]
Warum die frühere Version $64\pi^4$ falsch war, ist geometrisch:
$64\pi^4 / 16\pi^3 = 4\pi = S^2$ — sie zählte \textbf{eine 2-Sphäre zu
viel}. Dass der Fehler exakt eine ganze Sphären-Oberfläche beträgt,
bestätigt die Geometriezählung.

\section{Die wiederholte $\mu_0/\varepsilon_0$-Struktur [K]}

In den Ursprungsdokumenten (Dok.~011, 013, 015) sind $\mu_0$ und
$\varepsilon_0$ \emph{keine} unabhängigen Konstanten — sie folgen aus
$\alpha, e, c$:
\[
  \varepsilon_0 = \frac{e^2}{2\alpha h c}, \quad
  \mu_0 = \frac{2\alpha h}{e^2 c}, \quad
  \varepsilon_0\mu_0 c^2 = 1.
\]
Sie sind der Ort, wo Geometrie und Einheiten getrennt sichtbar werden:
$\mu_0 = \underbrace{4\pi}_{\text{Geometrie }S^2} \times
\underbrace{10^{-7}}_{\text{Einheit}}$. In der $\alpha$-Formel kürzt
sich das $4\pi$ exakt weg:
\[
  \alpha = \frac{e^2\mu_0 c}{4\pi\hbar}
         = \frac{e^2 \times 10^{-7} \times c}{\hbar}.
\]
Geometrie gegen Geometrie. In natürlichen Einheiten ($\hbar = c = 1$)
verschwinden $\mu_0,\varepsilon_0$ vollständig; $\alpha$ und $\xi$
bleiben als Zahlen. Setzt man zusätzlich $\alpha = 1$ (T0-Einheiten,
Dok.~011), wird die Ladung umdefiniert — $\alpha = e^2 = 1 \Rightarrow
e = 1$ — und auch $\alpha$ verschwindet. Dann bleibt \textbf{nur $\xi$}
als einziger freier dimensionsloser Parameter.

\textbf{Fazit:} Nichts wurde ans Ziel angepasst. Die $\pi$-Faktoren sind
durchgehend Sphären-Oberflächen der Feldgeometrie; eine saubere
Einheiten- und Dimensionszählung liefert sie zwingend.

\chapter{Warum QM und RT zu vermengen berechtigt ist}

Der übliche Einwand — „du mischst $\hbar$ (Quanten) mit $c$
(Relativität)" — wird nicht verteidigt, sondern umgekehrt: die
\textbf{Trennung} QM$|$RT ist das Künstliche, nicht die Vermengung.

\textbf{(1)} $\hbar$ und $c$ sind beide nur Einheiten-Brücken (Raum$\leftrightarrow$Zeit,
Energie$\leftrightarrow$Zeit). Keine ist eine Theorie. Beide gleichzeitig zu
benutzen ist erlaubt.

\textbf{(2)} Die T$^4$-Struktur braucht beide (P34): $c$ ist die
Raum-Zeit-Brücke, $\hbar$ die Masse-Zeit-Brücke. Setzt man $c=\hbar=1$,
kollabiert die Geometrie. Um sie zu sehen, muss man vermengen.

\textbf{(3)} $\tilde T \cdot m = 1$ \emph{ist} die Vermengung: Eigenzeit
(RT) mal Masse (QM). Die Physik vermengt ohnehin längst — Dirac, QFT,
Unruh, Hawking. FFGFT macht es nur explizit.

\textbf{(4)} Der Konflikt ist ein Interpretations-, kein
Formalismuskonflikt. Nach Dok.~230 ist die Instantaneität ein
Beschreibungs-Artefakt; auf T$^4$ sind Bell-Korrelationen topologisch
verbunden, nicht ontologisch nichtlokal. Korrigiert man die
Interpretation, ist die Vermengung notwendig.

\medskip
\textbf{Kernsatz.} Wer die Vermengung verbietet, verbietet die
Geometrie.

\chapter{Der Kernsatz}

\begin{quote}
Das Übergeordnete ist die Geometrie, die Resonanzen besitzt. Ihre
dimensionslosen Aussagen — $Q = 2/3$, die Generationsleiter (Exponenten
$3/2, 1, 2/3$), die harmonischen Vorfaktoren — legen alle
Verhältnisse fest, ohne einen einzigen Messwert. Der Messwert (eine
Masse in SI) ist nur der Kammerton: er sagt, \emph{wo} die Struktur
klingt, nicht \emph{wie} sie klingt.
\end{quote}

\chapter{Bezüge}

\textbf{Interne Dokumente:} Dok.~011 (Feinstruktur, $\mu_0/\varepsilon_0$),
Dok.~013 (SI, $\mu_0/\varepsilon_0$ aus $\alpha/e/c$), Dok.~015
(Natürliche Einheiten), Dok.~049 (Higgs-Integration), Dok.~060
(musikalische Resonanz), Dok.~158 (Koide-zu-g-2-Brücke), Dok.~190
(P33, P41–P43, Massenformel 190.3, Higgs-Formel 190.1), Dok.~230
(Bell auf T$^4$), Dok.~306 ($T\cdot E = 1$), Dok.~309 (Skalenanker).

\medskip
\noindent\textbf{Prüfskripte} (keine freien Parameter, alle Eingaben
deklariert): \texttt{ffgft\_310\_koide\_referenz.py},
\texttt{ffgft\_310\_massenleiter.py},
\texttt{ffgft\_310\_pi\_geometrie.py},
\texttt{ffgft\_310\_higgs\_zugang.py}.
